\documentclass[11pt]{article}
\usepackage[a4paper, margin=1in]{geometry}
\usepackage[T1]{fontenc}
\usepackage[utf8]{inputenc}
\usepackage{microtype}
\usepackage{parskip}
\usepackage{hyperref}
\hypersetup{colorlinks=true, urlcolor=blue}
\pagestyle{empty}

\begin{document}

\noindent
Claes Johnson, Professor emeritus of Applied Mathematics \\
KTH Royal Institute of Technology, Stockholm \quad\textbar\quad \texttt{claesjohnson@gmail.com}

\vspace{1em}

\noindent
To Prof.\ Louis H.\ Kauffman, Editor, \emph{Series on Knots and Everything} \\
World Scientific Publishing \quad\textbar\quad \texttt{kauffman@uic.edu} \hfill \today

\vspace{1em}

\noindent Dear Professor Kauffman,

\medskip

I should like to propose a monograph for the \emph{Series on Knots and
Everything}: \textbf{\emph{Real Quantum Mechanics: A Multiphase 3D Continuum
Formulation, Realised through Mind--AI Cooperation}} --- a complete 339-page
draft, Volume VI of the \emph{Applied Mathematics: Body and Soul} programme
(Volumes I--IV published by Springer, 2003--2009).

\medskip
\noindent\textbf{The book.}
The book develops a reformulation of quantum mechanics, \emph{RealQM}, that
returns to the real-three-dimensional-space picture Schr\"odinger originally
envisaged in 1926, before the wave function was moved onto $3N$-dimensional
configuration space. The $N$-electron wave function is a sum of one-electron
wave functions on non-overlapping spatial domains separated by free (Bernoulli)
boundaries; ground states minimise the ordinary Coulomb energy. There is no
configuration space, no probabilistic interpretation, and no antisymmetry
postulate --- Pauli exclusion is realised geometrically as spatial non-overlap.
The formulation is parameter-free at the atomic level and is validated across
chemistry (atoms, covalent bonding, hydrides, reactions, condensed phases),
extended to a one-paragraph proof of stability of matter and, in its most
recent development, to the atomic nucleus itself. There the same Coulomb
variational principle --- with protons and electrons simply exchanging their
inner and outer roles --- reproduces nuclear binding energies across the
periodic table: the alpha particle at $103\%$ of experiment and the
binding-per-nucleon curve, saturation and all, at $107$--$108\%$, from a single
scale fixed on the deuteron and with no strong force. The nuclear chapters
carry this into a Coulomb reading of fusion (a topological rearrangement of
proton and electron domains, including a neutrino-free route to solar hydrogen
burning) and into a direct comparison with quantum chromodynamics, which after
fifty-three years still yields no parameter-free prediction of any nuclear
binding energy. This nuclear material has grown into a small suite of companion
papers now under journal submission, and forms the most striking part of the
book.

\medskip
\noindent\textbf{Why this series.}
The book sits naturally in the \emph{Knots and Everything} tradition of
foundational, cross-disciplinary mathematical physics that the mainstream
journals and the large physics imprints tend to gate out. It is a constructive,
computational reformulation in the spirit of the series' established titles,
and it makes contact with topology (the free-boundary partition of space into
electron territories), variational analysis, and the philosophy of physics. It
is, frankly, a book that the conventional venues decline on editorial-policy
rather than technical grounds --- which is precisely why a series founded to
keep such work in print is its right home.

\medskip
\noindent\textbf{What is distinctive.}
Two features set the book apart from a conventional monograph. \emph{First}, it
is paired with a complete, runnable implementation: roughly 8000 lines of
JavaScript and WebGPU compute shaders that execute every quantitative claim
interactively in a web browser on a consumer laptop, curated as a public
Gallery. The book and the Gallery are designed to be read together; the reader
can change a parameter and watch the chemistry respond. \emph{Second}, the
implementation and a substantial part of the manuscript were produced through
extended collaboration with an AI assistant (Claude, Anthropic), and the book
documents that collaboration as part of its contribution --- a worked, honest
account of a new working method for mathematical science that will interest a
readership well beyond quantum chemistry.

\medskip
\noindent\textbf{Audience and credentials.}
The intended readership is graduate students and researchers in applied
mathematics, computational and theoretical chemistry, the foundations of
physics, and scientific computing, as well as readers following the
human--AI-collaboration question. The \emph{Body and Soul} series (with Johan
Hoffman and the FEniCS finite-element project) has a two-decade track record of
pairing mathematics with runnable computation; this volume carries that
programme to atomic-scale matter. (On a personal note relevant to the series:
I spent 1974--75 as a Dickson Instructor in the Mathematics Department at the
University of Chicago.)

The manuscript is complete and available immediately as a PDF, together with
the public code repository and Gallery. I should be glad to send the full draft
and to discuss the proposal further, and would welcome the series' technical
review.

\medskip

\noindent With my respectful regards,

\vspace{1.5em}

\noindent Claes Johnson \\
Professor emeritus of Applied Mathematics, KTH Royal Institute of Technology, Stockholm \\
\texttt{claesjohnson@gmail.com}

\end{document}
